\paragraph{Theo dõi}

\def\MOTA{\relax\ifmmode\mathrm{MOTA}\else MOTA\fi}
\def\MOTP{\relax\ifmmode\mathrm{MOTP}\else MOTP\fi}
\def\HOTA{\relax\ifmmode\mathrm{HOTA}\else HOTA\fi}
\def\IDTP{\relax\ifmmode\mathrm{IDTP}\else IDTP\fi}
\def\IDFN{\relax\ifmmode\mathrm{IDFN}\else IDFN\fi}
\def\IDFP{\relax\ifmmode\mathrm{IDFP}\else IDFP\fi}
\def\DetA{\relax\ifmmode\mathrm{DetA}\else DetA\fi}
\def\AssA{\relax\ifmmode\mathrm{AssA}\else AssA\fi}
\def\TPA{\relax\ifmmode\mathrm{TPA}\else TPA\fi}
\def\FNA{\relax\ifmmode\mathrm{FNA}\else FNA\fi}
\def\FPA{\relax\ifmmode\mathrm{FPA}\else FPA\fi}
\def\IDFone{\relax\ifmmode\mathrm{IDF}_1\else IDF1\fi}
\def\FP{\relax\ifmmode\mathrm{FP}\else FP\fi}
\def\FN{\relax\ifmmode\mathrm{FN}\else FN\fi}
\def\TP{\relax\ifmmode\mathrm{TP}\else TP\fi}

Để đánh giá chất lượng của phương pháp theo dõi đa đối tượng, các độ đo được sử
dụng bao gồm tập độ đo CLEAR (\MOTA, \MOTP), độ đo nhận dạng \IDFone, và độ đo
bậc cao \HOTA.

Tập độ đo CLEAR đánh giá hệ thống thông qua hai khía cạnh chính: độ chính xác
định vị và độ chính xác theo dõi. Độ chuẩn xác MOT (Multiple Object Tracking
Precision, \MOTP) thể hiện khả năng ước lượng vị trí chính xác của đối tượng
độc lập với khả năng nhận dạng. \MOTP\ được tính bằng trung bình sai số vị trí
của các cặp đối tượng được khớp thành công: \begin{equation}
  \MOTP = \frac{\sum_{i,t} d_t^i}{\sum_t c_t},
\end{equation} trong đó, $d_t^i$ là sai số khoảng cách giữa dự đoán và thực tế
của đối tượng thứ $i$ tại khung hình $t$, và $c_t$ là số lượng đối tượng được
khớp tại khung hình $t$.

Độ chính xác MOT (Multiple Object Tracking Accuracy, \MOTA) đánh giá độ chính
xác tổng thể trong việc phát hiện và duy trì quỹ đạo của đối tượng. \MOTA\ tổng
hợp ba loại lỗi: bỏ sót (misses), nhận diện sai (false positives) và sai số
ghép cặp/nhảy ID (mismatches): \begin{equation}
  \MOTA = 1 - \frac{\sum_t (m_t + fp_t + mme_t)}{\sum_t g_t},
\end{equation} trong đó, $m_t$ là số lượng bỏ sót, $fp_t$ là số nhận diện sai,
$mme_t$ là số lần nhảy ID, và $g_t$ là tổng số đối tượng thực tế tại khung hình
$t$.

Trong khi MOTA chú trọng nhiều vào độ chính xác phát hiện ở từng khung hình rời
rạc, \IDFone\ (Identification F1-score) tập trung đánh giá khả năng duy trì
danh tính (ID) cố định của đối tượng qua thời gian dài. \IDFone\ tính toán tỷ
lệ các phát hiện được gán đúng danh tính trên tổng số trung bình của các phát
hiện thực tế và dự đoán: \begin{equation}
  \IDFone = \frac{2 \times \IDTP}{2 \cdot \IDTP + \IDFP + \IDFN},
\end{equation} trong đó, \IDTP\ (Identification True Positives) là số lượng hộp
bao dự đoán đúng cả vị trí và danh tính, \IDFP\ (Identification False
Positives) là số hộp bao dự đoán sai danh tính hoặc vị trí, và \IDFN\
(Identification False Negatives) là số hộp bao thực tế bị bỏ sót hoặc gán nhầm
danh tính.

\HOTA\ là một độ đo mới nhằm khắc phục hạn chế của cả \MOTA\ và \IDFone\ bằng
cách cân bằng một cách rõ ràng giữa độ chính xác phát hiện (Detection), liên
kết (Association) và định vị (Localization) thành một chỉ số thống nhất duy
nhất.  Tại một ngưỡng định vị $\alpha$, \HOTA\ được định nghĩa là trung bình
nhân của điểm phát hiện và điểm liên kết: \begin{align}
  \DetA_\alpha &= \frac{\lvert \TP \rvert}{ \lvert \TP \rvert + \lvert \FN
  \rvert + \lvert \FP \rvert }, \\
  \AssA_\alpha &= \frac{1}{\lvert \TP \rvert} \sum_{c \in \TP} \mathcal{A}(c),
  \\
  \mathcal{A}(c) &= \frac{\lvert \TPA(c) \rvert}{ \lvert \TPA(c) \rvert +
  \lvert \FNA(c) \rvert + \lvert \FPA(c) \rvert}, \\
  \HOTA_\alpha &= \sqrt{\frac{\sum_{c \in \TP} \mathcal{A}(c)}{\lvert
  \TP \rvert + \lvert \FN \rvert + \lvert \FP \rvert}} = \sqrt{\DetA_\alpha
  \cdot \AssA_\alpha},
\end{align} trong đó, $\TP, \FN, \FP$ lần lượt là số lượng True Positives,
False Negatives, và False Positives trong phát hiện được tính trên ngưỡng
$\alpha$; $\TPA(c), \FNA(c), \FPA(c)$ lần lượt là số lượng phát hiện được gán
đúng ID của $c$, số lượng phát hiện là $c$ nhưng bị gán sai ID hoặc bị bỏ sót,
và số lượng phát hiện bị gán sai thành ID của $c$. Hàm $\mathcal{A}(c)$ thể
hiện độ chính xác liên kết (Association Accuracy) đối với mỗi True Positive
$c$.

Để phản ánh chính xác cả khả năng định vị, điểm \HOTA\ tổng quát được tính bằng
cách lấy tích phân (thường được xấp xỉ bằng giá trị trung bình) trên nhiều
ngưỡng định vị $\alpha$ khác nhau (ví dụ từ 0.05 đến 0.95):
\begin{equation}
    \HOTA = \int_0^1 \HOTA_\alpha\, d\alpha \approx \frac{1}{19} \sum_{\alpha \in
    \{0.05, 0.1, \dots, 0.95\}} \HOTA_\alpha.
\end{equation} Tương tự ta có thể định nghĩa \DetA\ và \AssA. \DetA\ thể hiện
độ chính xác phát hiện và \AssA\ thể hiện khả năng định danh nhất quán của bộ
theo dõi. Do đó, \HOTA\ có sự cân bằng đánh giá giữa phát hiện và liên kết.
